\documentclass{beamer}

\usepackage{graphicx}
\usepackage[style=authoryear]{biblatex}
\renewcommand*{\nameyeardelim}{\addcomma\addspace}

\addbibresource{../references.bib}
%\usepackage{epstopdf}

% Import the UNCMath theme.
% UNC is the default color theme.
\usepackage{source_files_slides/beamerthemeUNCMath}

% Uncomment to use a different color scheme.
%\definecolor{darkpurple}{rgb}{0.2,0,0.2}
%\setbeamercolor*{structure}{fg=black,bg=darkpurple}

\logo{

\includegraphics[height=3em]{STOR_712_Project/source_files_slides/oldwell_cmyk-eps-converted-to.pdf}
}

% A macro for making the title frame.
% Removes the bottom bar and logo temporarily.
% If you don't want these in other frames, 
% you could try mimicking this.
\newcommand{\titleframe}{
{
\setbeamertemplate{headline}{}
\setbeamertemplate{footline}{}
\begin{frame}
\titlepage
\end{frame}
}
}

%%%%%%%%%%%%%%%%%%%%%%%%%%%%%%%%%%%%%%%%%%
%
% Information for the title frame.
%
\title{Can We Improve \textbf{Personalized Federated Learning} With \textbf{Learning Rate Scheduling}?}
\subtitle{}
\author{Joe Lavond and Peter Lin}
\date{March 27, 2023}
%%%%%%%%%%%%%%%%%%%%%%%%%%%%%%%%%%%%%%%%%%
%
% Start of document.
%
\begin{document}

% This is the titleframe, do it this way 
% so that you don't have the logo on the title page.
\titleframe


\begin{frame}{Introduction to Federated Learning}
    \begin{definition}[Federated Learning (FL)]
        An iterative procedure allowing \textcolor{blue}{multiple users} to build a shared robust machine learning model \textcolor{red}{without sharing data} (e.g. healthcare). 
    \end{definition}

    Key pieces:
    \begin{itemize}
        \item Data: Multiple users with their data
        \item Training: Users each have a local copy of the model for training
        \item \textbf{Communication}: Send model information rather than data \\
            \quad (privacy) \\
            %\quad (privacy and resource demands) \\
    \end{itemize}

\end{frame}


\begin{frame}{Training}
    FL is an iterative process of training rounds. See Figure~\ref{fig: fl}.
    
    \begin{figure}
        \centering
        \includegraphics[width = .8\textwidth]{STOR_712_Project/images/fedavg.png}
        \caption{Single Iteration of Federated Training}
        \label{fig: fl}
    \end{figure}
    %
    The most common aggregation scheme for Step 4 in Figure~\ref{fig: fl} is FedAvg \parencite{fedavg}, which averages returning models.
\end{frame}


\begin{frame}{Personalized Federated Learning (pFL)}

\textcolor{red}{Q}: Is a single model appropriate if users have very different data sets? \\

\smallskip
\textcolor{blue}{A}: \textcite{pfl-bench} argues for individualized models tailored to each user's data.

\bigskip
How to personalize?
\begin{itemize}
    \item Main works Ditto \parencite{ditto}, FedBN \parencite{fedbn}, and FedEM \parencite{fedem} require additional computation and memory
    \item Simplest solution is fine-tuning from the single Federated Learning model
\end{itemize}
    
Is it possible to find a single model that, after fine-tuning, performs similarly to the more complex solutions?
\end{frame}


\begin{frame}{User Local Training}

    Within a round of FL, perform $K$ steps of GD on the current global model
    \begin{align}
        x^{j+1}=x^j-\textcolor{red}{\alpha}\nabla f(x^j)
    \end{align}

    Motivated by Taylor Expansion of \textcite{reptile}, we propose step size decay (within round)
    \begin{align}
        & x^{j+1}=x^j-\textcolor{blue}{\alpha\beta^{j-1}}\nabla f(x^j) \label{eq: lr-decay} \\
        & x^K = x^0 - \alpha\sum_{j=1}^K \beta^{j-1}\nabla f(x^{j-1}) \nonumber
    \end{align}

    Equations~\ref{eq: lr-decay} is a trade-off between accuracy ($\beta = 0$) and ability to fine-tune ($\beta = 1$). \\
    
\end{frame}


\begin{frame}[allowframebreaks]{Proposed Work}

    Theory: Establish convergence of FedAvg with our proposed learning rate.

    \bigskip
    Numerical: How does our learning rate scheduling impact model performance?
    \begin{enumerate}
        \item Split user data into train, validation, and test
        \item Hold out certain users to measure generalization
        \item Report average, standard deviation, and bottom ten percentile of model accuracy
    \end{enumerate}

    Compare with fixed learning rate FedAvg and Personalized Federated Learning methods.
    
\end{frame}

\begin{frame}[allowframebreaks]{References}
    \printbibliography
\end{frame}



\end{document}
